\section{Experimental Design}
\label{sec:experiments}
\begin{table}[!t]
\caption{Registered experimental settings}
\label{tab:settings}
\centering
\footnotesize
\begin{tabular}{@{}p{0.40\columnwidth}p{0.53\columnwidth}@{}}
\toprule
Setting & Value \\
\midrule
Learning/validation candidates & 60 / 15 instances \\
Feasible demonstration instances & 26 training / 4 validation \\
Held-out instances & 45 small / 30 medium / 15 large \\
Tasks and nodes per instance & 50 / 150 / 500 by scale \\
Workflow families & Montage, CyberShake, Epigenomics, LIGO, SIPHT shapes \\
Task CPU/memory demand & One configured unit each \\
Max. deployed active copies & 2 \\
Detection/replanning delay & 1.0 / 0.5 s \\
Switch/synchronization delay & 0.5 / 0.5 s \\
Link-recovery wait & 2.0 s \\
GCN layers/hidden width & 2 / 64 \\
Domain embedding / neighbors & 8 dimensions / 8 neighbors \\
GCN dropout / learning rate & 0.3 / $10^{-3}$ \\
Adam weight decay & $5\times10^{-4}$ \\
Negatives per positive & 5 \\
Pairwise ranking loss weight & 0.1 \\
Maximum epochs / patience & 100 / 10 \\
ALNS restart budget & 2 \\
Insertion GCN weight & 0.5 \\
Main small/medium budget & 1000 iterations, 10 seeds \\
Main large budget & 800 iterations, 5 seeds \\
\bottomrule
\end{tabular}
\end{table}

\begin{table}[!t]
\caption{SLA settings relative to the nominal HEFT-style reference}
\label{tab:sla}
\centering
\begin{tabular}{lcc}
\toprule
Scenario & $D/T_{\mathrm{ref}}$ & $R_{\mathrm{SLA}}/T_{\mathrm{ref}}$ \\
\midrule
Tight & 1.15 & 0.05 \\
Medium & 1.30 & 0.10 \\
Loose & 1.50 & 0.20 \\
\bottomrule
\end{tabular}
\end{table}

\subsection{Research Questions}
The independent evaluation addresses three questions.
\begin{itemize}
 \item \textbf{RQ1:} How do feasibility, feasible-solution cost, and recovery
 margin compare with nonlearned ALNS and the baseline suite?
 \item \textbf{RQ2:} Which search mechanisms and protection-policy choices
 account for observed changes, and are those changes supported by paired
 analysis?
 \item \textbf{RQ3:} How does learned guidance affect solution quality under
 different iteration budgets, and how does iteration efficiency differ from
 total elapsed time?
\end{itemize}

\subsection{Synthetic Instances and SLA Calibration}
Instances combine five custom DAG-shape generators named after Montage,
CyberShake, Epigenomics, LIGO, and SIPHT workflow families. Runtimes follow a
lognormal distribution with log-space mean 2 and standard deviation 0.7,
while dependency data volumes are sampled from a
configured range of 5--200 MB. These generators approximate broad structural
patterns; they are not real application executions or outputs of the
WfCommons generator. The current synthetic path assigns every task one unit
of CPU demand and one unit of memory demand. Workload heterogeneity therefore
comes primarily from runtime and dependency structure, while resource
capacity varies across nodes.

Synthetic machines draw CPU capacities from $\{8,16,32,64,96\}$ and memory
capacities from $\{16,32,64,128,256\}$. The medium-load setting samples
utilization between 0.30 and 0.60. Quality-based tier assignment combines
capacity, utilization, and variation, followed by tier-dependent prices and
failure-rate parameters. Links use tier-pair latency/bandwidth tables with
deterministic multiplicative perturbations. Fault domains are synthetic
hash-based groups, not observed racks or power domains. Resource data are
cluster-trace-inspired rather than measurements extracted from a production
trace.

For each instance, a passive HEFT-style reference defines
$T_{\mathrm{ref}}$. Deadline and recovery window are scaled as shown in
Table~\ref{tab:sla}. Small, medium, and large instances contain 50, 150, and
500 tasks and the same number of resource nodes, respectively. The held-out
set contains 45, 30, and 15 instances across these scales, for 90 instances
in total. Each SLA setting has 30 test instances across all scales. We do not
retune the SLA parameters after observing the test-set feasibility rates.

\subsection{Independent Learning and Model Selection}
The input protocol fixes 60 learning candidates and 15 validation candidates
before teacher generation. Workflow-content and resource-snapshot hashes
prevent cross-role reuse, including renamed or SLA-modified copies. The
teacher runs only on these inputs and retains an instance only when it finds
at least one verifier-feasible demonstration. In the completed rebuild, 26
training and four validation instances contributed demonstrations. Failed
teacher instances remain assigned to their original roles; they are not
replaced by solver-test instances.

Crucially, supervised-label availability does not redefine the validation or
test population. All 15 original validation instances are retained for
end-to-end solver comparison, including those for which the teacher found
no feasible demonstration. The 90 test instances use separate random-seed
blocks and are never supplied to the teacher or training loss.

Candidate training configurations use the same frozen teacher dataset.
Validation loss controls early stopping.
End-to-end validation uses five repeated solver seeds per instance, with
1000 ALNS iterations as the maximum budget. The registered model-selection
rule first prefers more feasible validation runs. If feasibility ties, it
compares relative cost on identically paired runs where both candidates are
feasible; a remaining tie follows the preregistered candidate priority.
All other search
parameters, including the insertion weight of 0.5, remain fixed.

Validation selected the placement-only checkpoint, with the heuristic
protection rule described in Section~\ref{sec:method}. This single checkpoint
is frozen before testing and used throughout the main, component, and budget
studies. Alternative checkpoints and their diagnostic comparisons are
archived as development records and do not redefine the final model using
test-set outcomes.

\subsection{Baseline Configurations}
\label{sec:baselines}
The suite contains 14 solver labels. Eight built-in labels represent
HEFT-style placement, PEFT-style placement, no-fault-aware placement,
reactive-only placement, full active replication, critical-path replication,
criticality-based hybrid protection, and reliability-aware placement.
HEFT and No-FT share the same passive placement logic in this implementation
and should not be counted as independent algorithmic evidence. Every
returned deployment is evaluated with the same nominal scheduler, recovery
policy, and cost model.

Four additional configurations are external references:
\begin{itemize}
 \item \textbf{RECR-D adaptation:} Retains resource penalties, distributed
 sub-deadlines, and bounded replica promotion. Energy terms are removed;
 placement is performed in the evaluator's canonical order to avoid a
 mismatch between planned and replayed schedules. Recovery-budget feedback
 is mapped to the common simulator rather than the source reliability model.
 \item \textbf{HFWPF adaptation:} Retains the slack-based choice between
 replication and reexecution and a capacity-biased placement rule. The
 source method's execution-time uncertainty, elastic virtual machines, and
 online adaptation are not reproduced.
 \item \textbf{GA-Placement:} Uses a conventional genetic search over placement
 and protection choices as a generic evolutionary reference. It is not a
 reproduction of a genetic-programming hyper-heuristic.
 \item \textbf{RA-MOMA adaptation:} Uses a multi-objective memetic search with
 HEFT-style initialization and problem-specific neighborhoods. The resource
 model is fixed, order genes are omitted, and source reliability and billing
 terms are mapped to the evaluator's cost and recovery-risk objectives.
\end{itemize}
These are explicitly static adaptations, not bit-exact reproductions of
the cited systems \cite{taghinezhad2025recrd,ren2026hfwpf,qin2023ramoma}.
Results under this common gate assess the adapted configurations on
RLC-HFTSP; they do not rank the source papers on their original tasks.
Nonlearned ALNS uses the same search framework as GCN-ALNS, with all learned
hooks disabled.

\subsection{Registered Studies and Computational Budgets}
The main comparison uses ten repeated seeds for stochastic methods on small
and medium instances and five on large instances. The maximum ALNS budgets
are 1000 and 800 iterations, respectively. Deterministic references run once
per instance. The main comparison therefore contains 4200 planned records.
The repeated runs measure optimizer variability on fixed instances; they are
not additional independent workflow samples. Different algorithm families
retain their registered internal search parameters, so the main study is
not a matched-wall-clock comparison across all methods.

The component study uses small and medium test instances, ten repeated
seeds, and 1000 iterations. Its 14 configurations include the complete solver
and removals of learned guidance, warm start, learned replica ranking,
learned destruction, feasibility-first scoring, restarts, initialization
diversity, recovery penalty, deployment risk, domain separation, dynamic
rescheduling, and two repair-mode preference variants. This produces
10\,500 planned records.

These changes belong to different categories. Disabling dynamic rescheduling
restricts reexecution to the task's original domains and changes the recovery
policy; it does not eliminate recovery. Removing domain separation relaxes
an admissibility constraint and is not a pure search-operator ablation.
The variants originally labelled all-active and all-passive apply a mode
preference when repairing removed tasks. Initialization and other repair
actions can retain or change modes, so these configurations do not impose a
strict homogeneous policy on every task. Their results describe repair-mode
preferences. A clean comparison of strictly all-active, all-passive, and
hybrid deployments requires globally constrained, separately versioned runs.
These changes must be distinguished from removing learned ranking or a search
restart.

The budget study compares ALNS and GCN-ALNS on small instances at 50, 100,
200, 500, and 1000 iterations, with ten seeds each, yielding 4500 records.
These three studies contain 19\,200 planned records. Their scopes and budgets
are fixed before inspecting the complete test results. Additional
training-configuration diagnostics are archived separately; they are not an
independent research question or a basis for replacing the frozen checkpoint.

\subsection{Metrics and Statistical Analysis}
Primary metrics are SLA-feasible rate, feasible-deployment cost,
$M(\chi)=\reff-\rmax$, and the iteration at which the solver first records a
verified feasible deployment. We also record full solver runtime, final
verification time, one-time learned inference time, replica usage, and
scenario-set size. Reported verification time for an individual result is
the final evaluation call; it does not account for every verification call
inside ALNS. It cannot alone identify the full solver's dominant runtime
component.

Feasible rate is first averaged over repeated seeds for each instance, then
macro-averaged over instances. For the main descriptive cost table, feasible
seeds are first averaged within each instance, then instances with at least
one feasible run receive equal weight. These method-specific cost means are
descriptive because their successful instance sets can differ. Comparative
cost, recovery margin, and first-feasible iteration are paired by the same
instance and seed and restricted to runs where both methods succeed; unknown
first-feasible iteration markers are excluded. Pairwise differences are then
aggregated within instances so the workflow instance, rather than every
optimizer seed, is the independent statistical unit.

We report paired mean differences, 95\% instance-bootstrap intervals with
2000 resamples, paired $t$-tests, and Wilcoxon signed-rank tests where defined.
For the component comparisons, Holm adjustment is performed separately for
each reported metric family. A feasibility increase, a lower cost among
different successful subsets, and a higher conditional safety margin are
not interchangeable findings. Likewise, fewer iterations to feasibility
must not be presented as a directly measured reduction in wall-clock
time-to-feasibility.

\subsection{Execution Environment and Reproducibility}
The current run uses Windows 11, Python 3.12.4, NumPy 1.26.4, and PyTorch
2.5.1 with CUDA 12.1 on one NVIDIA A40 GPU. The host exposes 32 logical CPU
processors. GPU execution is used for training; per-instance score inference
in the scheduling pipeline is performed on CPU. Independent solver runs use
eight worker processes for ordinary stages and 12 for the large main study,
with per-worker numerical-library threads limited to one.

The insertion weights in~\eqref{eq:insertion} are
$(1,0.5,1,2,0.5,0.5)$; deadline and recovery penalties are 200 and 100.
ALNS uses initial temperature 1, cooling factor 0.995, a 50-iteration
adaptation segment, reaction factor 0.1, maximum removal fraction 0.2,
and a 300-iteration stagnation threshold with at most two restarts.
The recovery-node weights in~\eqref{eq:recovery_selection} are
$(1,1,1,0.5,0.2)$. Communication, passive reserve, and recovery-proxy
weights are 0.01, 0.1, and 1, respectively. Tier base failure-rate
parameters are $10^{-6}$, $5\times10^{-6}$, and $10^{-5}$ for cloud,
edge, and device nodes, with configured load/variation adjustments.
Costs are normalized model units rather than measured currency charges.

Each experiment stores configuration, source-code, model, and instance
fingerprints together with the returned deployment and protection decisions.
Resume checks reject incompatible versions. The source snapshot includes
the correction that serializes reexecutions on recovery nodes and enforces
CPU/memory footprint checks at the final gate. Earlier results produced by
the previous verifier are excluded from the independent comparisons.
